\section{Separable algebraic extensions}
\label{section-separable-extensions}

\noindent
In characteristic $p$ something funny happens with irreducible polynomials
over fields. We explain this in the following lemma.

\begin{lemma}
\label{lemma-irreducible-polynomials}
Let $F$ be a field. Let $P \in F[x]$ be an irreducible polynomial over $F$.
Let $P' = \text{d}P/\text{d}x$ be the derivative of $P$ with respect
to $x$. Then one of the following two cases happens
\begin{enumerate}
\item $P$ and $P'$ are relatively prime, or
\item $P'$ is the zero polynomial.
\end{enumerate}
The second case can only happen if $F$ has characteristic $p > 0$.
In this case $P(x) = Q(x^q)$ where $q = p^f$ is a power of $p$ and
$Q \in F[x]$ is an irreducible polynomial such that $Q$ and $Q'$
are relatively prime.
\end{lemma}

\begin{proof}
Note that $P'$ has degree $< \deg(P)$. Hence if $P$ and $P'$ are not relatively
prime, then $(P, P') = (R)$ where $R$ is a polynomial of degree $< \deg(P)$
contradicting the irreducibility of $P$. This proves we have the dichotomy
between (1) and (2).

\medskip\noindent
Assume we are in case (2) and $P = a_d x^d + \ldots + a_0$. Then
$P' = da_d x^{d - 1} + \ldots + a_1$. In characteristic $0$ we see
that this forces $a_d, \ldots, a_1 = 0$ which would mean $P$ is constant
a contradiction. Thus we conclude that the characteristic $p$ is positive.
In this case the condition $P' = 0$ forces $a_i = 0$ whenever $p$ does
not divide $i$.
In other words, $P(x) = P_1(x^p)$ for some nonconstant polynomial $P_1$.
Clearly, $P_1$ is irreducible as well. By induction on the degree we
see that $P_1(x) = Q(x^q)$ as in the statement of the lemma, hence
$P(x) = Q(x^{pq})$ and the lemma is proved.
\end{proof}

\begin{definition}
\label{definition-separable}
Let $F$ be a field. Let $K/F$ be an extension of fields.
\begin{enumerate}
\item We say an irreducible polynomial $P$ over $F$ is {\it separable}
if it is relatively prime to its derivative.
\item Given $\alpha \in K$ algebraic over $F$ we say $\alpha$ is
{\it separable} over $F$ if its minimal polynomial is separable over $F$.
\item If $K$ is an algebraic extension of $F$, we say $K$ is
{\it separable}\footnote{For nonalgebraic extensions
this definition does not make sense and is not the correct one. We refer
the reader to Algebra, Sections \ref{algebra-section-separability} and
\ref{algebra-section-separability-continued}.}
over $F$ if every element of $K$ is separable over $F$.
\end{enumerate}
\end{definition}

\noindent
By Lemma \ref{lemma-irreducible-polynomials} in characteristic $0$ every
irreducible polynomial is separable, every algebraic element in an extension
is separable, and every algebraic extension is separable.

\begin{lemma}
\label{lemma-separable-goes-up}
Let $K/E/F$ be a tower of algebraic field extensions.
\begin{enumerate}
\item If $\alpha \in K$ is separable over $F$, then $\alpha$ is separable
over $E$.
\item if $K$ is separable over $F$, then $K$ is separable over $E$.
\end{enumerate}
\end{lemma}

\begin{proof}
We will use Lemma \ref{lemma-irreducible-polynomials} without further mention.
Let $P$ be the minimal polynomial of $\alpha$ over $F$.
Let $Q$ be the minimal polynomial of $\alpha$ over $E$.
Then $Q$ divides $P$ in the polynomial ring $E[x]$, say $P = QR$.
Then $P' = Q'R + QR'$. Thus if $Q' = 0$, then $Q$ divides $P$ and $P'$
hence $P' = 0$ by the lemma. This proves (1). Part (2)
follows immediately from (1) and the definitions.
\end{proof}

\begin{lemma}
\label{lemma-recognize-separable}
Let $F$ be a field. An irreducible polynomial $P$ over $F$
is separable if and only if $P$ has pairwise distinct roots in an
algebraic closure of $F$.
\end{lemma}

\begin{proof}
Suppose that $\alpha \in \overline{F}$ is a root of both $P$ and $P'$.
Then $P = (x - \alpha)Q$ for some polynomial $Q$. Taking derivatives
we obtain $P' = Q + (x - \alpha)Q'$. Thus $\alpha$ is a root of $Q$.
Hence we see that if $P$ and $P'$ have a common root, then $P$
does not have pairwise distinct roots. Conversely, if $P$ has
a repeated root, i.e., $(x - \alpha)^2$ divides $P$, then $\alpha$
is a root of both $P$ and $P'$. Combined with
Lemma \ref{lemma-relatively-prime-polynomials} this proves the lemma.
\end{proof}

\begin{lemma}
\label{lemma-nr-roots-unchanged}
Let $F$ be a field and let $\overline{F}$ be an algebraic closure of $F$.
Let $p > 0$ be the characteristic of $F$. Let $P$ be a polynomial
over $F$. Then the set of roots of $P$ and $P(x^p)$ in $\overline{F}$
have the same cardinality (not counting multiplicity).
\end{lemma}

\begin{proof}
Clearly, $\alpha$ is a root of $P(x^p)$ if and only if $\alpha^p$ is a
root of $P$. In other words, the roots of $P(x^p)$ are the roots of
$x^p - \beta$, where $\beta$ is a root of $P$. Thus it suffices to show
that the map $\overline{F} \to \overline{F}$, $\alpha \mapsto \alpha^p$
is bijective. It is surjective, as $\overline{F}$ is algebraically closed
which means that every element has a $p$th root. It is injective because
$\alpha^p = \beta^p$ implies $(\alpha - \beta)^p = 0$ because
the characteristic is $p$. And of course in a field $x^p = 0$ implies
$x = 0$.
\end{proof}

\noindent
Let $F$ be a field and let $P$ be an irreducible polynomial over $F$.
Then we know that $P = Q(x^q)$ for some separable irreducible polynomial $Q$
(Lemma \ref{lemma-irreducible-polynomials}) where $q$ is a power of
the characteristic $p$ (and if the characteristic is zero, then
$q = 1$\footnote{A good convention for this chapter is to set $0^0 = 1$.}
and $Q = P$). By Lemma \ref{lemma-nr-roots-unchanged} the number of
roots of $P$ and $Q$ in any algebraic closure of $F$ is the same.
By Lemma \ref{lemma-recognize-separable} this number is equal to the degree
of $Q$.

\begin{definition}
\label{definition-separable-degree}
Let $F$ be a field. Let $P$ be an irreducible polynomial over $F$.
The {\it separable degree} of $P$ is the cardinality of the
set of roots of $P$ in any algebraic closure of $F$ (see discussion
above). Notation $\deg_s(P)$.
\end{definition}

\noindent
The separable degree of $P$ always divides the degree and the quotient
is a power of the characteristic. If the characteristic is zero, then
$\deg_s(P) = \deg(P)$.

\begin{situation}
\label{situation-finitely-generated}
Here $F$ be a field and $K/F$ is a finite extension generated by elements
$\alpha_1, \ldots, \alpha_n \in K$. We set $K_0 = F$ and
$$
K_i = F(\alpha_1, \ldots, \alpha_i)
$$
to obtain a tower of finite extensions
$K = K_n / K_{n - 1} / \ldots / K_0 = F$.
Denote $P_i$ the minimal polynomial of $\alpha_i$ over $K_{i - 1}$.
Finally, we fix an algebraic closure $\overline{F}$ of $F$.
\end{situation}

\noindent
Let $F$, $K$, $\alpha_i$, and $\overline{F}$ be as in
Situation \ref{situation-finitely-generated}.
Suppose that $\varphi : K \to \overline{F}$ is a morphism of extensions
of $F$. Then we obtain maps $\varphi_i : K_i \to \overline{F}$.
The image of $P_i \in K_{i - 1}[x]$ by $\varphi_{i - 1}$
is a polynomial $P_i^{\varphi_{i - 1}} \in \overline{F}[x]$
and $\varphi(\alpha_i)$ is a root of this polynomial.

\begin{lemma}
\label{lemma-count-embeddings}
In Situation \ref{situation-finitely-generated} the correspondence
$$
\Mor_F(K, \overline{F})
\longrightarrow
\{(\beta_1, \ldots, \beta_n)\text{ as below}\},
\quad
\varphi \longmapsto (\varphi(\alpha_1), \ldots, \varphi(\alpha_n))
$$
is a bijection. Here the right hand side is the set of $n$-tuples
$(\beta_1, \ldots, \beta_n)$ of elements of $\overline{F}$ having
the following property:
\begin{enumerate}
\item $\beta_1 \in \overline{F}$ is a root of $P_1$;
let $\varphi_1 : K_1 \to \overline{F}$ be the homomorphism
over $F$ sending $\alpha_1$ to $\beta_1$,
\item $\beta_2 \in \overline{F}$ is a root of $P_2^{\varphi_1}$;
let $\varphi_2 : K_2 \to \overline{F}$ be the homomorphism
extending $\varphi_1$ sending $\alpha_2$ to $\beta_2$,
\item and so on until,
\item $\beta_n \in \overline{F}$ is a root of $P_n^{\varphi_{n - 1}}$.
\end{enumerate}
In each step the homorphism $\varphi_i$ exists and is unique because
$K_i = K_{i - 1}[x]/(P_i)$ and $\beta_i$ is a root of $P_i^{\varphi_{i - 1}}$.
\end{lemma}

\begin{proof}
The map from left to right is discussed above the lemma.
Let $(\beta_1, \ldots, \beta_n)$ be an element of the right hand side.
Then we let $\varphi : K = K_n \to \overline{F}$ be the unique homorphism
extending $\varphi_{n - 1}$ sending $\alpha_n$ to $\beta_n$.
Uniqueness implies that the two constructions are mutually inverse.
\end{proof}

\begin{lemma}
\label{lemma-count-embeddings-explicitly}
In Situation \ref{situation-finitely-generated} we have
$|\Mor_F(K, \overline{F})| = \prod_{i = 1}^n \deg_s(P_i)$.
\end{lemma}

\begin{proof}
This follows immediately from Lemma \ref{lemma-count-embeddings}.
Observe that a key ingredient we are tacitly using here is the
well-definedness of the separable degree of an irreducible polynomial
which was observed just prior to
Definition \ref{definition-separable-degree}.
\end{proof}

\noindent
We now use the result above to characterize separable field extensions.

\begin{lemma}
\label{lemma-separably-generated-separable}
Assumptions and notation as in Situation \ref{situation-finitely-generated}.
If each $P_i$ is separable, i.e., each $\alpha_i$ is separable over
$K_{i - 1}$, then
$$
|\Mor_F(K, \overline{F})| = [K : F]
$$
and the field extension $K/F$ is separable. If one of the $\alpha_i$ is
not separable over $K_{i - 1}$, then
$|\Mor_F(K, \overline{F})| < [K : F]$.
\end{lemma}

\begin{proof}
If $\alpha_i$ is separable over $K_{i - 1}$ then
$\deg_s(P_i) = \deg(P_i) = [K_i : K_{i - 1}]$
(last equality by Lemma \ref{lemma-degree-minimal-polynomial}).
By multiplicativity (Lemma \ref{lemma-multiplicativity-degrees}) we have
$$
[K : F] = \prod [K_i : K_{i - 1}] = \prod \deg(P_i) =
\prod \deg_s(P_i) = |\Mor_F(K, \overline{F})|
$$
where the last equality is Lemma \ref{lemma-count-embeddings-explicitly}.
By the exact same argument we get the strict inequality
$|\Mor_F(K, \overline{F})| < [K : F]$ if one of the $\alpha_i$ is
not separable over $K_{i - 1}$.

\medskip\noindent
Finally, assume again that each $\alpha_i$ is separable over $K_{i - 1}$.
We will show $K/F$ is separable.
Let $\gamma = \gamma_1 \in K$ be arbitrary. Then we can find additional
elements $\gamma_2, \ldots, \gamma_m$ such that
$K = F(\gamma_1, \ldots, \gamma_m)$ (for example we could take
$\gamma_2 = \alpha_1, \ldots, \gamma_{n + 1} = \alpha_n$).
Then we see by the last part of the lemma (already proven above)
that if $\gamma$ is not separable over $F$ we would have the
strict inequality $|\Mor_F(K, \overline{F})| < [K : F]$
contradicting the very first part of the lemma (already prove above
as well).
\end{proof}

\begin{lemma}
\label{lemma-separable-equality}
Let $K/F$ be a finite extension of fields. Let $\overline{F}$ be an
algebraic closure of $F$. Then we have
$$
|\Mor_F(K, \overline{F})| \leq [K : F]
$$
with equality if and only if $K$ is separable over $F$.
\end{lemma}

\begin{proof}
This is a corollary of Lemma \ref{lemma-separably-generated-separable}.
Namely, since $K/F$ is finite we can find finitely many elements
$\alpha_1, \ldots, \alpha_n \in K$ generating $K$ over $F$ (for example
we can choose the $\alpha_i$ to be a basis of $K$ over $F$).
If $K/F$ is separable, then each $\alpha_i$ is separable over
$F(\alpha_1, \ldots, \alpha_{i - 1})$ by Lemma \ref{lemma-separable-goes-up}
and we get equality by Lemma \ref{lemma-separably-generated-separable}.
On the other hand, if we have equality, then no matter how we choose
$\alpha_1, \ldots, \alpha_n$ we get that $\alpha_1$ is separable over
$F$ by Lemma \ref{lemma-separably-generated-separable}. Since we
can start the sequence with an arbitrary element of $K$ it follows
that $K$ is separable over $F$.
\end{proof}

\begin{lemma}
\label{lemma-separable-permanence}
Let $E/k$ and $F/E$ be separable algebraic extensions of fields. Then $F/k$
is a separable extension of fields.
\end{lemma}

\begin{proof}
Choose $\alpha \in F$. Then $\alpha$ is separable algebraic over $E$.
Let $P = x^d + \sum_{i < d} a_i x^i$ be the minimal polynomial of
$\alpha$ over $E$. Each $a_i$ is separable algebraic over $k$.
Consider the tower of fields
$$
k \subset k(a_0) \subset k(a_0, a_1) \subset \ldots \subset
k(a_0, \ldots, a_{d - 1}) \subset k(a_0, \ldots, a_{d - 1}, \alpha)
$$
Because $a_i$ is separable algebraic over $k$ it is separable algebraic
over $k(a_0, \ldots, a_{i - 1})$ by Lemma \ref{lemma-separable-goes-up}.
Finally, $\alpha$ is separable algebraic over $k(a_0, \ldots, a_{d - 1})$
because it is a root of $P$ which is irreducible
(as it is irreducible over the possibly bigger field $E$)
and separable (as it is separable over $E$).
Thus $k(a_0, \ldots, a_{d - 1}, \alpha)$ is separable over $k$
by Lemma \ref{lemma-separably-generated-separable}
and we conclude that $\alpha$ is separable over $k$ as desired.
\end{proof}

\begin{lemma}
\label{lemma-separable-elements}
Let $E/k$ be a field extension. Then the elements of $E$ separable
over $k$ form a subextension of $E/k$.
\end{lemma}

\begin{proof}
Let $\alpha, \beta \in E$ be separable over $k$. Then $\beta$ is separable
over $k(\alpha)$ by Lemma \ref{lemma-separable-goes-up}.
By Lemma \ref{lemma-separably-generated-separable} (applied with $n = 2$,
$\alpha_1 = \alpha$, and $\alpha_2 = \beta$) we see that
$k(\alpha, \beta)$ is separable over $k$.
\end{proof}





